% concatenated from main.tex
\documentclass[]{article}

\usepackage{amsfonts}
\usepackage{amsmath}
\usepackage{amssymb}
\usepackage{amsthm}
\usepackage{algorithm}
\usepackage{algorithmicx}
\usepackage{algpseudocode}
\usepackage{graphicx}
\usepackage{indentfirst}
\usepackage{placeins}
\usepackage[hidelinks]{hyperref}
\usepackage{url}

\renewcommand{\baselinestretch}{1.2}
\topmargin 0cm
\oddsidemargin 0.28cm
\evensidemargin 0.28cm
\textwidth 16cm
\textheight 22.8cm
\headheight 0cm
\headsep 0cm
\emergencystretch=2em

\numberwithin{equation}{section}
\theoremstyle{plain}
\newtheorem{theorem}{Theorem}[section]
\newtheorem{lemma}[theorem]{Lemma}
\newtheorem{proposition}[theorem]{Proposition}
\theoremstyle{remark}
\newtheorem{remark}[theorem]{Remark}
\numberwithin{algorithm}{section}
\algrenewcommand\algorithmicrequire{\textbf{Input:}}
\algrenewcommand\algorithmicensure{\textbf{Output:}}

\newcommand{\R}{\mathbb{R}}
\newcommand{\Mr}{\mathcal{M}_r}
\newcommand{\Mle}{\mathcal{M}_{\le r}}
\newcommand{\eps}{\varepsilon}
\newcommand{\rank}{\operatorname{rank}}
\newcommand{\range}{\operatorname{range}}

\hypersetup{
 pdftitle={Robust High-Order Projector-Splitting Integrators},
 pdfauthor={Shiheng Zhang and Jingwei Hu},
 pdfkeywords={dynamical low-rank approximation; projector splitting; exactness; numerical increments; robust integrators; small singular values; Runge--Kutta methods}
}

\begin{document}

\title{Robust High-Order Projector-Splitting Integrators}
\author{Shiheng Zhang\thanks{Department of Applied Mathematics,
University of Washington, Seattle, WA 98195, USA
(\texttt{shzhang3@uw.edu}).}
\and
Jingwei Hu\thanks{Department of Applied Mathematics,
University of Washington, Seattle, WA 98195, USA
(\texttt{hujw@uw.edu}).}}
\date{}

\maketitle

\begin{abstract}
We develop a general framework for constructing robust high-order projector-splitting
integrators for dynamical low-rank approximation. For a prescribed matrix increment, we show that the standard \(K\)-\(S\)-\(L\) projector-splitting step \cite{LubichOseledets2014} is equivalent to a reduced \(K\)-\(L\) step, thereby eliminating the explicit
backward \(S\)-step. We then establish a central relaxed exactness property of the standard
projector-splitting integrator: the rank-$r$ approximation inherits the accuracy of the numerical matrix increment, with an error bound independent of small singular values. The resulting schemes evolve fixed rank-\(r\) factors and require neither basis augmentation nor rank truncation.
As concrete examples, we combine the framework with selected second- and third-order Runge--Kutta methods to obtain robust high-order projector-splitting integrators.
Numerical experiments confirm the predicted uniform convergence rates with respect to small singular values.
\end{abstract}

\medskip
\noindent\textbf{Keywords.}
Dynamical low-rank approximation; projector splitting; exactness; numerical
increment; robust low-rank integrator; small singular values; Runge--Kutta
method.

\smallskip
\noindent\textbf{MSC 2020.} 65L05, 65L20, 65F55.

\section{Introduction}
\label{sec:introduction}

Dynamical low-rank approximation of time-dependent matrices
\cite{KochLubich2007} replaces the solution
\(A(t)\in\R^{m\times n}\) of a large matrix differential equation
\begin{equation}\label{eq:full-ode}
 \dot A(t)=F(t,A(t)),\qquad A(t_0)=A_0,
\end{equation}
by matrices \(Y(t)\) of rank \(r\), represented in the factorized form
\begin{equation}\label{eq:factorization}
 Y=USV^\top,
\end{equation}
where \(U\in\R^{m\times r}\) and \(V\in\R^{n\times r}\) have
orthonormal columns and \(S\in\R^{r\times r}\) is a small matrix.  The
rank constraint is imposed by projecting the right-hand side of
\eqref{eq:full-ode} onto the tangent space of the rank-\(r\) manifold,
\begin{equation}\label{eq:projected-ode}
 \dot Y(t)=P_r(Y(t))F(t,Y(t)).
\end{equation}
Throughout the paper, \(\|\cdot\|\) denotes the Frobenius norm and
\(\|\cdot\|_2\) denotes the matrix operator norm.

The differential equations for the factors in \eqref{eq:factorization}
contain \(S^{-1}\) in the equations for \(U\) and \(V\).  Their direct
integration by standard methods therefore becomes problematic when \(S\)
has small singular values.  Such values occur naturally when the prescribed
rank is larger than the effective rank or when the singular values of the
solution have no distinct gap.  An error estimate is called robust when its
constants and stepsize threshold are independent of small singular values
of the low-rank approximations.

The projector-splitting integrator of Lubich and Oseledets
\cite{LubichOseledets2014} avoids the inverse of the small factor \(S\) by
integrating the three terms of the tangent space projector in turn.
For a known matrix curve, these subproblems can be integrated exactly,
and the resulting \(K\)-\(S\)-\(L\) step depends only on the matrix
increment \(\Delta A\).  We use a numerical approximation of this
increment to construct higher-order methods.

For any prescribed increment, the \(K\)-\(S\)-\(L\) step reduces to a
reduced \(K\)-\(L\) step.  With \(V_0\) an orthonormal basis for the row space
at the beginning of the step, the two factorizations and the output are
\begin{equation}\label{eq:two-factorizations-intro}
 (Y_0+\Delta A)V_0=U_1\widehat S_1,
 \qquad (Y_0+\Delta A)^\top U_1=V_1S_1^\top,
 \qquad Y_1=U_1S_1V_1^\top.
\end{equation}
The algebraic simplification of the core update eliminates the backward
\(S\)-step.  The reduced step requires only products of \(\Delta A\)
and \(\Delta A^\top\) with rank-\(r\) factors.

The projector-splitting integrator has an exactness property
\cite{LubichOseledets2014}.  Suppose that a rank-\(r\) matrix
\(Y_\star\) satisfies the row-space overlap condition
\(Y_\star=(Y_\star V_0)Q\) for some \(Q\in\R^{r\times n}\).
If \(Y_0+\Delta A=Y_\star\), the step gives \(Y_1=Y_\star\).

We prove a relaxed exactness property for approximate increments.  If \(Y_1\)
is the rank-admissible output for a prescribed \(\Delta A\), and the
comparison matrix satisfies \(Y_\star=(Y_\star V_0)Q\), then
\begin{equation}\label{eq:stable-exactness-intro}
 \|Y_1-Y_\star\|\le(1+\|Q\|_2)\|Y_0+\Delta A-Y_\star\|,
\end{equation}
with a constant determined by the overlap matrix \(Q\).  In particular,
\(Y_0+\Delta A=Y_\star\) gives \(Y_1=Y_\star\).

To compare the approximation with the exact solution, we construct a
rank-\(r\) curve satisfying the overlap condition with a uniformly bounded
matrix \(Q\).  Its distance to the exact solution is bounded in terms of
the normal component of \(F\).
For a computation at \(t_n:=t_0+nh\), let \(\Delta A_n\) be the prescribed
increment and let \(A_n(s)\) solve
\(\dot A_n(s)=F(t_n+s,A_n(s))\), \(A_n(0)=Y_n\).
Under the standing assumptions on \(F\), we obtain
\begin{equation}\label{eq:stable-transfer-intro}
 \|Y_{n+1}-A_n(h)\|
 \le C\bigl(\|Y_n+\Delta A_n-A_n(h)\|+h\eps_r\bigr),
\end{equation}
where \(\eps_r\) bounds the normal component and \(C\) is independent of the
singular values of the numerical approximations.

Set
\begin{equation}\label{eq:increment-defect-intro}
 d_h:=\max_{t_{n+1}\le T}\|Y_n+\Delta A_n-A_n(h)\|.
\end{equation}
For every rank-admissible computation, in which the matrices to be
factorized have rank \(r\), the global error satisfies
\begin{equation}\label{eq:general-global-intro}
 \max_{t_n\le T}\|Y_n-A(t_n)\|
 \le C_T\left(\delta+\eps_r+\frac{d_h}{h}\right),
\end{equation}
with \(\delta\) bounding the initial error.  The increment rule enters this
bound only through \(d_h\).

As an application, we use the reduced \(K\)-\(L\) step to compute the stages of
Runge--Kutta (RK) methods.  At each stage, the increment is a linear combination
of previously computed values of \(F\).  The reduced \(K\)-\(L\) step starts
from \((U_n,S_n,V_n)\), and \(F\) is then evaluated at the new stage
approximation.  Every stage therefore uses the row basis \(V_n\), so
\eqref{eq:stable-transfer-intro} applies throughout the calculation.  The explicit
midpoint rule and Heun's three-stage third-order rule then give
\(d_h=O(h^3+h^2\eps_r)\) and \(d_h=O(h^4+h^2\eps_r)\), and hence the global
bound \(C_T(\delta+\eps_r+h^p)\) with \(p=2\) and \(p=3\).
Both methods keep rank \(r\) throughout the computation and require
neither basis augmentation nor rank truncation.

Robust first-order error bounds for projector splitting were proved in
\cite{KieriLubichWalach2016}.  Projected
Runge--Kutta methods of orders two and three use retraction and rank
truncation at the stages \cite{KieriVandereycken2019}.  Other robust
constructions include the basis-update-and-Galerkin (BUG) methods
\cite{CerutiLubich2022}, a rank-adaptive integrator
\cite{CerutiKuschLubich2022}, the midpoint BUG method
\cite{CerutiEinkemmerKuschLubich2024}, parallel low-rank integrators
\cite{Kusch2025}, and high-order Runge--Kutta BUG methods
\cite{NobileRiffaud2026}.  The XL integrator uses a \(K\)-step to update
one basis, augments that basis, and then performs an \(L\)-step
\cite{einkemmer2025asymptotic}.

Runge--Kutta approximations of the matrix increment have also been used in
a single projector-splitting step \cite{HochbruckNeherSchrammer2023}.
Here the reduced \(K\)-\(L\) step is applied at every Runge--Kutta stage.
The \(K\)-\(S\)-\(L\) step also has an extended-retraction interpretation
\cite{AbsilOseledets2015,SeguinCerutiKressner2024}.
The counterexample in \cite{Zhang2026Counterexample} shows that the
state-dependent Strang composition can have a nonzero \(O(h^2)\) local
error under the standard robust assumptions.  The methods studied here
use the exactness property for prescribed increments to obtain higher
robust orders.

Section~\ref{sec:psi} derives the reduced \(K\)-\(L\) step and recalls
the exactness property.  Sections~\ref{sec:increments}--\ref{sec:global}
give the assumptions and error analysis for general numerical increments.
Section~\ref{sec:rk} constructs the Runge--Kutta methods and proves their
orders, and Section~\ref{sec:numerical} presents numerical experiments.
The comparison curve is constructed in Appendix~\ref{app:comparison-curve}.

\section{The projector-splitting integrator}
\label{sec:psi}

\subsection{The splitting}
\label{sec:splitting}

At \(Y=USV^\top\) of rank \(r\), the orthogonal projector onto the tangent
space of the rank-\(r\) manifold, applied to \(Z\in\R^{m\times n}\), is
\begin{equation}\label{eq:tangent-projector}
 P_r(Y)Z=ZVV^\top-UU^\top ZVV^\top+UU^\top Z.
\end{equation}
To derive the projector-splitting integrator, suppose first that the
matrix curve \(A(t)\) is known, and consider
\begin{equation}\label{eq:prescribed-curve-dlra}
 \dot Y(t)=P_r(Y(t))\dot A(t).
\end{equation}
The three terms of \eqref{eq:tangent-projector} are integrated in turn
\cite{LubichOseledets2014}.  Set \(t_1:=t_0+h\).  Over \([t_0,t_1]\),
starting from \(Y_0=U_0S_0V_0^\top\), the factor equations are
\begin{equation}\label{eq:subflows}
 \begin{aligned}
  \dot K(t)&=\dot A(t)V_0,
  &&K(t_0)=U_0S_0,\\
  \dot S(t)&=-U_1^\top\dot A(t)V_0,
  &&S(t_0)=\widehat S_1,\\
  \dot L(t)&=\dot A(t)^\top U_1,
  &&L(t_0)=V_0\widetilde S_0^\top.
 \end{aligned}
\end{equation}
Between the substeps one factorizes \(K(t_1)=U_1\widehat S_1\) and
\(L(t_1)=V_1S_1^\top\), with \(U_1\) and \(V_1\) having orthonormal
columns, and sets \(\widetilde S_0:=S(t_1)\).  The step returns
\(Y_1=U_1S_1V_1^\top\).  These equations avoid \(S^{-1}\).
The middle equation is the backward \(S\)-step arising from the negative
term in \eqref{eq:tangent-projector}.

\subsection{The \texorpdfstring{\(K\)-\(S\)-\(L\)}{K-S-L} step for a prescribed increment}
\label{sec:practical-step}

Set
\begin{equation}\label{eq:exact-matrix-increment}
 \Delta A:=A(t_1)-A(t_0)=\int_{t_0}^{t_1}\dot A(t)\,dt.
\end{equation}
Integrating \eqref{eq:subflows} gives the following algorithm
\cite{LubichOseledets2014}.  It depends on the prescribed curve only through
\(\Delta A\), so the same formulas apply to any given matrix
\(\Delta A\in\R^{m\times n}\).  Starting from
\begin{equation}\label{eq:ksl-start}
 Y_0=U_0S_0V_0^\top,\qquad U_0^\top U_0=V_0^\top V_0=I_r,
 \qquad S_0\ \hbox{invertible},
\end{equation}
the \(K\)-\(S\)-\(L\) step consists of the following calculations.

\begin{enumerate}
\item Set
\begin{equation}\label{eq:ksl-k}
 K_1:=U_0S_0+\Delta A\,V_0.
\end{equation}
Compute a factorization
\begin{equation}\label{eq:ksl-k-factor}
 U_1\widehat S_1=K_1,
\end{equation}
where \(U_1\) has orthonormal columns and \(\widehat S_1\) is invertible.

\item Set
\begin{equation}\label{eq:ksl-s}
 \widetilde S_0:=\widehat S_1-U_1^\top\Delta A\,V_0.
\end{equation}

\item Set
\begin{equation}\label{eq:ksl-l}
 L_1:=V_0\widetilde S_0^\top+\Delta A^\top U_1.
\end{equation}
Compute a factorization
\begin{equation}\label{eq:ksl-l-factor}
 V_1S_1^\top=L_1,
\end{equation}
where \(V_1\) has orthonormal columns and \(S_1\) is invertible.
\end{enumerate}
The output is
\begin{equation}\label{eq:ksl-output}
 Y_1:=U_1S_1V_1^\top.
\end{equation}

\subsection{The reduced \texorpdfstring{\(K\)-\(L\)}{K-L} step}
\label{sec:two-factorizations}

For a prescribed increment, the middle substep can be eliminated
algebraically.  Equations \eqref{eq:ksl-k}--\eqref{eq:ksl-s} give
\begin{equation}\label{eq:ksl-s-simplified}
 \widehat S_1=U_1^\top U_0S_0+U_1^\top\Delta A\,V_0,
 \qquad \widetilde S_0=U_1^\top U_0S_0.
\end{equation}
Substitution in \eqref{eq:ksl-l} yields
\begin{equation}\label{eq:ksl-l-simplified}
 L_1=V_0S_0^\top U_0^\top U_1+\Delta A^\top U_1
     =(Y_0+\Delta A)^\top U_1.
\end{equation}

Consequently, the following reduced \(K\)-\(L\) step gives the same matrix
\(Y_1\):
\begin{enumerate}
\item Set
\begin{equation}\label{eq:two-factor-k}
 K_1:=U_0S_0+\Delta A\,V_0=(Y_0+\Delta A)V_0.
\end{equation}
Compute the factorization
\begin{equation}\label{eq:two-factor-k-factor}
 U_1\widehat S_1=K_1,
\end{equation}
where \(U_1\) has orthonormal columns.

\item Set
\begin{equation}\label{eq:two-factor-l}
 L_1:=V_0S_0^\top U_0^\top U_1+\Delta A^\top U_1
     =(Y_0+\Delta A)^\top U_1.
\end{equation}
Compute the factorization
\begin{equation}\label{eq:two-factor-l-factor}
 V_1S_1^\top=L_1,
\end{equation}
where \(V_1\) has orthonormal columns, and set
\(Y_1:=U_1S_1V_1^\top\).
\end{enumerate}
The calculation is \emph{rank-admissible} when \(\rank K_1=r\).
Indeed,
\begin{equation}\label{eq:second-factor-rank}
 L_1^\top V_0
 =U_1^\top(Y_0+\Delta A)V_0
 =U_1^\top K_1
 =\widehat S_1
 \quad\Longrightarrow\quad \rank L_1=r.
\end{equation}
A computation with several steps or stages is rank-admissible when this
condition holds at every reduced \(K\)-\(L\) step.  Only the products of \(\Delta A\)
and \(\Delta A^\top\) with the rank-\(r\) factors are required.

The two factorizations give
\begin{equation}\label{eq:ksl-old-row-identities}
 Y_1=U_1L_1^\top=U_1U_1^\top(Y_0+\Delta A),
 \qquad U_1U_1^\top(Y_0+\Delta A)V_0=(Y_0+\Delta A)V_0.
\end{equation}

\subsection{Exactness}
\label{sec:exactness}

The exactness property of \cite[Theorem~4.1]{LubichOseledets2014} follows
from \eqref{eq:ksl-old-row-identities}.  It uses the row-space overlap
condition
\begin{equation}\label{eq:old-row-overlap}
 Y_\star=(Y_\star V_0)Q
 \qquad\text{for some }Q\in\R^{r\times n}.
\end{equation}
For a rank-\(r\) matrix \(Y_\star=U_\star S_\star V_\star^\top\), condition
\eqref{eq:old-row-overlap} holds exactly when \(V_\star^\top V_0\) is
invertible, which is the standard overlap condition between the new and old
row spaces.  The matrix \(Q\) is then unique,
\begin{equation}\label{eq:overlap-matrix}
 Q=(V_\star^\top V_0)^{-1}V_\star^\top,
 \qquad
 \|Q\|_2=\frac{1}{\sigma_{\min}(V_\star^\top V_0)},
\end{equation}
and \(\sigma_{\min}(V_\star^\top V_0)\) is the cosine of the largest
principal angle between \(\range(V_\star)\) and \(\range(V_0)\).  Thus
\(\|Q\|_2\) measures how much the two row spaces overlap.

\begin{proposition}
\label{prop:exactness}
Let \(Y_1\) be the rank-admissible output obtained from
\eqref{eq:ksl-start} with prescribed increment \(\Delta A\).  If
\(Y_0+\Delta A\) satisfies \eqref{eq:old-row-overlap}, then
\[
 Y_1=Y_0+\Delta A.
\]
\end{proposition}

\begin{proof}
Write \(Y_\star:=Y_0+\Delta A\) and let \(Q\) be as in
\eqref{eq:old-row-overlap}.  The second identity in
\eqref{eq:ksl-old-row-identities} gives
\((I-U_1U_1^\top)Y_\star V_0=0\), and therefore
\[
 (I-U_1U_1^\top)Y_\star=(I-U_1U_1^\top)(Y_\star V_0)Q=0.
\]
The first identity in \eqref{eq:ksl-old-row-identities} then gives
\(Y_1=U_1U_1^\top Y_\star=Y_\star\).
\end{proof}

\section{Numerical increments}
\label{sec:increments}

A numerical approximation \(\Delta A_n\) to the solution increment
defines a time-stepping method through the reduced \(K\)-\(L\) step.  At step
\(n\), apply \eqref{eq:two-factor-k}--\eqref{eq:two-factor-l-factor} with
\((U_0,S_0,V_0,\Delta A)\) replaced by
\((U_n,S_n,V_n,\Delta A_n)\).
The first-order method in \cite{LubichOseledets2014} uses the Euler
increment \(\Delta A_n=hF(t_n,Y_n)\).  Other numerical approximations of
the increment can be used in the same way.  For the Runge--Kutta methods
below, the reduced \(K\)-\(L\) step also computes the internal stage approximations.

\subsection{Standing assumptions}
\label{sec:assumptions}

Fix \(T>0\) and \(1\le r<\min(m,n)\), and set
\begin{equation}\label{eq:rank-strata}
 \Mr:=\{Y\in\R^{m\times n}:\rank(Y)=r\}.
\end{equation}
The following assumptions hold for all \(0\le t\le T\) and all matrices
of the indicated sizes.

\begin{enumerate}
\item The map \(F:[0,T]\times\R^{m\times n}\to\R^{m\times n}\) is
continuous, uniformly bounded, and globally Lipschitz-continuous in its
matrix argument:
\begin{equation}\label{eq:assumption-BL}
 \|F(t,X)\|\le B,
 \qquad
 \|F(t,X)-F(t,Z)\|\le L_F\|X-Z\|.
\end{equation}

\item For some \(\eps_r\ge0\), the normal component satisfies, for every
\(Y=USV^\top\in\Mr\),
\begin{equation}\label{eq:assumption-normal}
 \|(I-UU^\top)F(t,Y)(I-VV^\top)\|\le\eps_r.
\end{equation}
\end{enumerate}

\subsection{The error in the numerical increment}
\label{sec:defect}

Let \(t_n=t_0+nh\), and suppose that the approximation at step \(n\) is
stored as
\begin{equation}\label{eq:method-start}
 Y_n=U_nS_nV_n^\top,
 \qquad U_n^\top U_n=V_n^\top V_n=I_r,
 \qquad S_n\ \hbox{invertible}.
\end{equation}
At every step, prescribe a matrix
\(\Delta A_n\in\R^{m\times n}\), apply the reduced \(K\)-\(L\) step to
\((U_n,S_n,V_n)\), and denote its output by
\(Y_{n+1}\).  The increment may depend on previously computed stages and
steps.  Let \(A_n(s)\) solve
\begin{equation}\label{eq:increment-local-solution}
 \dot A_n(s)=F(t_n+s,A_n(s)),\qquad A_n(0)=Y_n,
\end{equation}
and define
\begin{equation}\label{eq:increment-defect}
 d_n:=\|Y_n+\Delta A_n-A_n(h)\|,
 \qquad d_h:=\max_{t_{n+1}\le T}d_n.
\end{equation}
Here \(A_n(s)\) is the exact solution at elapsed time \(s\), starting
from \(Y_n\).  Thus \(d_n\) is the error in the increment used at step \(n\).

\section{Relaxed exactness}
\label{sec:stable-exactness}

We first extend Proposition~\ref{prop:exactness} to approximate increments.
The resulting estimate compares the reduced \(K\)-\(L\) approximation \(Y_1\)
with any matrix \(Y_\star\) satisfying \eqref{eq:old-row-overlap}.

\begin{theorem}
\label{thm:stable-exactness}
Let \(Y_1\) be the rank-admissible output obtained from
\eqref{eq:ksl-start} with prescribed increment \(\Delta A\), and let
\(Y_\star\in\R^{m\times n}\) satisfy \eqref{eq:old-row-overlap}.  Then
\begin{equation}\label{eq:ksl-stability}
 \|Y_1-Y_\star\|\le(1+\|Q\|_2)\|Y_0+\Delta A-Y_\star\|.
\end{equation}
\end{theorem}

\begin{proof}
By \eqref{eq:ksl-old-row-identities},
\begin{align}
 Y_1-Y_\star
 &=U_1U_1^\top(Y_0+\Delta A-Y_\star)
   -(I-U_1U_1^\top)Y_\star,                             \label{eq:ksl-proof-1}\\
 (I-U_1U_1^\top)Y_\star
 &=(I-U_1U_1^\top)Y_\star V_0Q
  =(I-U_1U_1^\top)(Y_\star-Y_0-\Delta A)V_0Q.           \label{eq:ksl-proof-2}
\end{align}
Since \(U_1U_1^\top\) is an orthogonal projector and \(\|V_0\|_2=1\),
\eqref{eq:ksl-proof-1}--\eqref{eq:ksl-proof-2} imply
\eqref{eq:ksl-stability}.
\end{proof}

For the one-step estimates, fix \(t_0\) and \(Y_0\), and let \(A(s)\)
denote the exact solution of
\begin{equation}\label{eq:local-exact-solution}
 \dot A(s)=F(t_0+s,A(s)),\qquad A(0)=Y_0,
 \qquad 0\le s\le T-t_0.
\end{equation}
We apply \eqref{eq:ksl-stability} with
\(Y_\star=\widetilde Y(s)\), where the following lemma gives a
rank-\(r\) comparison curve and a uniform bound on \(\|Q_s\|_2\).

\begin{lemma}
\label{lem:comparison-curve}
Under the standing assumptions, there
exists \(h_*>0\), depending only on \(m,n,r\), and \(L_F\), with the following
property.  Let \(Y_0\in\Mr\), let \(V_0\) have orthonormal columns with
\(Y_0=Y_0V_0V_0^\top\), and let \(A\) solve
\eqref{eq:local-exact-solution}.  For \(0\le t_0\le T-h\) and
\(0<h\le h_*\), there is an
absolutely continuous curve \(\widetilde Y:[0,h]\to\Mr\) satisfying
\begin{equation}\label{eq:comparison-rank}
 \widetilde Y(0)=Y_0,
 \qquad \rank\bigl(\widetilde Y(s)\bigr)=r,
\end{equation}
and, for almost every \(0\le s\le h\),
\begin{equation}\label{eq:comparison-residual}
 \|\dot{\widetilde Y}(s)-F(t_0+s,\widetilde Y(s))\|
 \le\sqrt{mn}\,\eps_r.
\end{equation}
Moreover,
\begin{align}
 \|\widetilde Y(s)-A(s)\|
 &\le\sqrt{mn}\,\eps_r s e^{L_Fs},                    \label{eq:comparison-flow}\\
 \widetilde Y(s)&=(\widetilde Y(s)V_0)Q_s,
 \qquad Q_s\in\R^{r\times n},\quad \|Q_s\|_2\le3.   \label{eq:comparison-overlap}
\end{align}
All these bounds are independent of the singular values of \(Y_0\).
\end{lemma}

The proof is given in Appendix~\ref{app:comparison-curve}.

Combining Theorem~\ref{thm:stable-exactness} and
Lemma~\ref{lem:comparison-curve}
gives the following one-step error estimate.

\begin{theorem}
\label{thm:accuracy-transfer}
Under the standing assumptions, let
\(Y_0=U_0S_0V_0^\top\in\Mr\) and let \(A\) solve
\eqref{eq:local-exact-solution}.  For \(0\le t_0\le T-s\) and
\(0\le s\le h_*\), let \(Y_1\) be the rank-admissible output for a
prescribed increment \(\Delta A\).  Then
\begin{equation}\label{eq:accuracy-transfer}
 \|Y_1-A(s)\|
 \le4\|Y_0+\Delta A-A(s)\|
   +5\sqrt{mn}\,\eps_r s e^{L_Fs}.
\end{equation}
\end{theorem}

\begin{proof}
Apply Lemma~\ref{lem:comparison-curve} to \(Y_0\) and \(V_0\).
Equations \eqref{eq:ksl-stability} and
\eqref{eq:comparison-overlap} give
\[
 \|Y_1-\widetilde Y(s)\|
 \le4\|Y_0+\Delta A-\widetilde Y(s)\|.
\]
Hence
\begin{align*}
 \|Y_1-A(s)\|
 &\le\|Y_1-\widetilde Y(s)\|+\|\widetilde Y(s)-A(s)\|\\
 &\le4\|Y_0+\Delta A-A(s)\|+5\|\widetilde Y(s)-A(s)\|\\
 &\le4\|Y_0+\Delta A-A(s)\|+5\sqrt{mn}\,\eps_r s e^{L_Fs}.
\end{align*}
\end{proof}

In particular,
\begin{equation}\label{eq:accuracy-transfer-consequence}
 \|Y_0+\Delta A-A(s)\|=O(h^q)
 \quad\Longrightarrow\quad
 \|Y_1-A(s)\|=O(h^q+h\eps_r)
\end{equation}
for \(0\le s\le h\le h_*\).  If \(\eps_r=0\) and the prescribed increment is
exact, \eqref{eq:accuracy-transfer} reduces to
Proposition~\ref{prop:exactness}.

At step \(n\), the theorem applies with \(Y_0,\Delta A,Y_1,A(s)\)
replaced by \(Y_n,\Delta A_n,Y_{n+1},A_n(s)\), respectively.

\section{Global error analysis}
\label{sec:global}

We use Theorem~\ref{thm:accuracy-transfer} to bound the global error in
terms of the increment error \(d_h\) defined in \eqref{eq:increment-defect}.

\begin{theorem}
\label{thm:general-increments}
Under the standing assumptions, let
\(A(t)\) solve \eqref{eq:full-ode} on \([t_0,T]\).  Suppose that
\(Y_0=U_0S_0V_0^\top\in\Mr\), with \(U_0,V_0\) having orthonormal columns
and \(S_0\) invertible, satisfies
\begin{equation}\label{eq:initial-error}
 \|Y_0-A(t_0)\|\le\delta.
\end{equation}
There exist \(h_0>0\) and \(C_T>0\), depending only on
\(m,n,r,B,L_F\), and \(T\), such that every rank-admissible sequence generated
by prescribed increments \(\Delta A_n\), with \(0<h\le h_0\), satisfies
\begin{equation}\label{eq:general-increment-global-error}
 \max_{t_n\le T}\|Y_n-A(t_n)\|
 \le C_T\left(\delta+\eps_r+\frac{d_h}{h}\right).
\end{equation}
In particular, if \(p\ge1\) and \(C_{\rm inc}\) is independent of \(h\),
\begin{equation}\label{eq:increment-order-transfer}
 d_h\le C_{\rm inc}(h^{p+1}+h\eps_r)
 \quad\Longrightarrow\quad
 \max_{t_n\le T}\|Y_n-A(t_n)\|
 \le C_T'(\delta+\eps_r+h^p).
\end{equation}
Here \(C_T'\) may also depend on \(C_{\rm inc}\).  The constants are
independent of the singular values along the numerical trajectory.
\end{theorem}

\begin{proof}
Set \(e_n:=\|Y_n-A(t_n)\|\).  Theorem~\ref{thm:accuracy-transfer}, applied
to \(\Delta A_n\) with \(s=h\), and the Lipschitz stability of
\eqref{eq:full-ode} give
\begin{align}
 e_{n+1}
 &\le \|Y_{n+1}-A_n(h)\|+\|A_n(h)-A(t_{n+1})\| \notag\\
 &\le 4d_n+5\sqrt{mn}\,\eps_rh e^{L_Fh}+e^{L_Fh}e_n.
 \label{eq:general-increment-recurrence}
\end{align}
Consequently,
\begin{equation}\label{eq:general-increment-iteration}
 e_n\le e^{L_Fnh}\delta+
 \sum_{k=0}^{n-1}e^{L_F(n-1-k)h}
 \left(4d_k+5\sqrt{mn}\,\eps_rh e^{L_Fh}\right).
\end{equation}
For \(t_n\le T\), \(n\le T/h\) and every exponential in the sum is at
most \(e^{L_FT}\).  Choosing \(h_0\le\min(1,h_*)\) in
\eqref{eq:general-increment-iteration} proves
\eqref{eq:general-increment-global-error}.  Substitution of the hypothesis
in \eqref{eq:increment-order-transfer} proves its conclusion.
\end{proof}

\section{Runge--Kutta examples}
\label{sec:rk}

We now construct Runge--Kutta methods in which the reduced \(K\)-\(L\) step
computes each internal stage and the final approximation.  The
Runge--Kutta coefficients determine the increments, and \(F\) is
evaluated at the resulting low-rank stage approximations.
Every reduced \(K\)-\(L\) step starts from the factors \((U_n,S_n,V_n)\).

\subsection{The Runge--Kutta scheme and its error bound}
\label{sec:rk-criterion}

For an explicit Runge--Kutta method with \(s_{\mathrm{RK}}\) stages,
write its tableau as
\[
 \begin{array}{c|ccc}
  c_1&a_{11}&\cdots&a_{1s_{\mathrm{RK}}}\\
  \vdots&\vdots&&\vdots\\
  c_{s_{\mathrm{RK}}}&a_{s_{\mathrm{RK}}1}&\cdots&a_{s_{\mathrm{RK}}s_{\mathrm{RK}}}\\ \hline
   &b_1&\cdots&b_{s_{\mathrm{RK}}}
 \end{array},
 \qquad a_{ij}=0\ (j\ge i),\qquad
 c_i=\sum_{j<i}a_{ij}.
\]
Here \(c_i\in[0,1]\) specifies the stage time and \(b_i\) is its final
weight.  Starting with \(Y_{n,1}:=Y_n\), the stage increments and final
increment are
\begin{equation}\label{eq:common-base-increments}
 \begin{aligned}
  \Delta A_{n,i}
  &:=h\sum_{j=1}^{i-1}a_{ij}F(t_n+c_jh,Y_{n,j}),
  &&2\le i\le s_{\mathrm{RK}},\\
  \Delta A_n
  &:=h\sum_{j=1}^{s_{\mathrm{RK}}}b_jF(t_n+c_jh,Y_{n,j}).
 \end{aligned}
\end{equation}
For each \(i=2,\ldots,s_{\mathrm{RK}}\), apply the reduced \(K\)-\(L\) step to
\((U_n,S_n,V_n)\) with
\(\Delta A_{n,i}\), obtaining \(Y_{n,i}\), and then evaluate
\(F(t_n+c_i h,Y_{n,i})\).  After all stages, the same calculation with
\(\Delta A_n\) gives \(Y_{n+1}\).

\begin{algorithm}[htbp]
\caption{One step of the projector-splitting Runge--Kutta method.}
\label{alg:common-base-rk}
\begin{algorithmic}[1]
\Require \(t_n,h\), the factors in \eqref{eq:method-start}, and
the explicit tableau \((a_{ij},b_i,c_i)\)
\State Set \(Y_{n,1}:=Y_n\) and \(F_{n,1}:=F(t_n,Y_n)\).
\For{\(i=2,\ldots,s_{\mathrm{RK}}\)}
 \State Set \(\Delta A_{n,i}:=h\sum_{j=1}^{i-1}a_{ij}F_{n,j}\).
 \State Apply the reduced \(K\)-\(L\) step \eqref{eq:two-factor-k}--\eqref{eq:two-factor-l-factor}
 to \((U_n,S_n,V_n)\) with increment \(\Delta A_{n,i}\), and denote
 its output by \(Y_{n,i}\).
 \State Set \(F_{n,i}:=F(t_n+c_i h,Y_{n,i})\).
\EndFor
\State Set \(\Delta A_n:=h\sum_{i=1}^{s_{\mathrm{RK}}}b_iF_{n,i}\).
\State Apply the reduced \(K\)-\(L\) step to \((U_n,S_n,V_n)\) with increment
\(\Delta A_n\).
\Ensure Its output \(Y_{n+1}=U_{n+1}S_{n+1}V_{n+1}^\top\).
\end{algorithmic}
\end{algorithm}
\FloatBarrier

For the analysis, let \(A_n\) be the exact solution of
\eqref{eq:increment-local-solution} and set
\begin{equation}\label{eq:flow-field-defined}
 g_n(s):=F(t_n+s,A_n(s))=\dot A_n(s).
\end{equation}
Theorem~\ref{thm:accuracy-transfer} applies to every rank-admissible
stage in Algorithm~\ref{alg:common-base-rk}.  The following theorem
combines the stage estimates with the final Runge--Kutta weights.

\begin{theorem}
\label{thm:rk-stage-criterion}
Under the standing assumptions, fix a
positive integer \(p\) and an explicit tableau as in
\eqref{eq:common-base-increments}, with real coefficients independent of
\(h\), \(n\), and the singular values.  Suppose that
\begin{equation}\label{eq:rk-quadrature-moments}
 \sum_{i=1}^{s_{\mathrm{RK}}}b_i=1,\qquad
 \sum_{i=1}^{s_{\mathrm{RK}}}b_i c_i^k=\frac1{k+1}
 \quad(1\le k\le p-1).
\end{equation}
For \(0<h\le h_0\le\min(1,h_*)\) and \(0\le t_n\le T-h\), consider a
rank-admissible computation of Algorithm~\ref{alg:common-base-rk} from
\(Y_n\in\Mr\).  Assume that \(g_n\in C^p([0,h])\), with
\begin{equation}\label{eq:rk-criterion-assumptions}
 \sup_{0\le s\le h}\|g_n^{(p)}(s)\|\le M_p,\qquad
 \|Y_{n,i}-A_n(c_i h)\|\le C_{\rm st}(h^p+h\eps_r)
 \quad\text{if }b_i\ne0.
\end{equation}
The bounds \(M_p,C_{\rm st}\) hold uniformly for \(0<h\le h_0\), and
\(h_0,M_p,C_{\rm st}\) are independent of the singular values.  Then
\begin{equation}\label{eq:rk-criterion-local}
 d_n\le C(h^{p+1}+h^2\eps_r),\qquad
 \|Y_{n+1}-A_n(h)\|\le C(h^{p+1}+h\eps_r).
\end{equation}
If the same \(h_0,M_p,C_{\rm st}\) apply at every step, then for the
solution \(A(t)\) of \eqref{eq:full-ode} under the initial-error assumption
\eqref{eq:initial-error},
\[
 \max_{t_n\le T}\|Y_n-A(t_n)\|
 \le C_T(\delta+\eps_r+h^p).
\]
The constant \(C\) may depend on \(m,n,r,B,L_F,p\), the fixed tableau,
and \(M_p,C_{\rm st}\); \(C_T\) may also depend on \(T\).
Both are independent of the singular values of the numerical approximations.
\end{theorem}

\begin{proof}
Since \(A_n(h)=Y_n+\int_0^h g_n(s)\,ds\),
\begin{align}
 Y_n+\Delta A_n-A_n(h)
 &=h\sum_{i=1}^{s_{\mathrm{RK}}}b_i
   \bigl(F(t_n+c_i h,Y_{n,i})-g_n(c_i h)\bigr)\notag\\
 &\quad+h\sum_{i=1}^{s_{\mathrm{RK}}}b_i g_n(c_i h)
       -\int_0^h g_n(s)\,ds.
 \label{eq:rk-criterion-decomposition}
\end{align}
Taylor expansion of \(g_n\) through degree \(p-1\), followed by
\eqref{eq:rk-quadrature-moments}, gives
\begin{equation}\label{eq:rk-criterion-quadrature}
 \left\|h\sum_{i=1}^{s_{\mathrm{RK}}}b_i g_n(c_i h)
        -\int_0^h g_n(s)\,ds\right\|
 \le \frac{M_p}{p!}
 \left(\frac1{p+1}+\sum_{i=1}^{s_{\mathrm{RK}}}|b_i|c_i^p\right)h^{p+1}.
\end{equation}
Consequently,
\begin{align}
 d_n
 &\le C h^{p+1}
   +hL_F\sum_{i=1}^{s_{\mathrm{RK}}}|b_i|
           \|Y_{n,i}-A_n(c_i h)\|\notag\\
 &\le C h^{p+1}
   +L_F C_{\rm st}\sum_{i=1}^{s_{\mathrm{RK}}}|b_i|
        (h^{p+1}+h^2\eps_r)
 \le C(h^{p+1}+h^2\eps_r).
 \label{eq:rk-criterion-defect}
\end{align}
The local error estimate follows from Theorem~\ref{thm:accuracy-transfer}.
If the hypotheses hold uniformly over the steps, Theorem
\ref{thm:general-increments} gives the global estimate.
\end{proof}


For the two examples below, let \(p\in\{2,3\}\) and assume
\begin{equation}\label{eq:assumption-smoothness}
 \begin{aligned}
  F&\in C^p\bigl([0,T]\times\R^{m\times n};\R^{m\times n}\bigr),\\
  \|\partial_t^aD_X^bF(t,X)\|&\le M,
  \qquad a,b\in\mathbb N_0,\quad 1\le a+b\le p.
 \end{aligned}
\end{equation}
Here \(D_X^bF\) is the \(b\)th Fr\'echet derivative with respect to \(X\),
with the multilinear operator norm induced by the Frobenius norm.

\begin{lemma}
\label{lem:flow-derivatives}
Under the standing assumptions and
\eqref{eq:assumption-smoothness}, there are constants \(M_1,\ldots,M_p\),
depending only on \(B,L_F,M\), such that
\begin{equation}\label{eq:flow-derivative-bounds}
 \|g_n^{(q)}(s)\|\le M_q,
 \qquad 1\le q\le p,\quad 0\le s\le h.
\end{equation}
These bounds are uniform over all \(0\le t_n\le T-h\) and initial
matrices \(Y_n\).
\end{lemma}

\begin{proof}
All derivatives of \(F\) below are evaluated at
\((t_n+s,A_n(s))\).  For \(p\ge2\),
\begin{align}
 g_n'&=\partial_tF+D_XF[g_n],                         \label{eq:g-first}\\
 g_n''&=\partial_t^2F+2D_X\partial_tF[g_n]+D_X^2F[g_n,g_n]
          +D_XF[g_n'].                              \label{eq:g-second}
\end{align}
For \(p=3\), additionally,
\begin{align}
 g_n'''&=\partial_t^3F+3D_X\partial_t^2F[g_n]
 +3D_X^2\partial_tF[g_n,g_n]+D_X^3F[g_n,g_n,g_n]\notag\\
 &\quad+3D_X\partial_tF[g_n']+3D_X^2F[g_n,g_n']+D_XF[g_n''].
                                                     \label{eq:g-third}
\end{align}
Since \(\|g_n\|\le B\) and \(\|D_XF\|\le L_F\), the bounds follow
successively from these identities.
\end{proof}

\subsection{The explicit midpoint method}
\label{sec:midpoint}

We use the explicit midpoint Runge--Kutta method
\begin{equation}\label{eq:midpoint-tableau}
\begin{array}{c|cc}
0&0&0\\
\frac12&\frac12&0\\ \hline
&0&1
\end{array}.
\end{equation}

Starting from \eqref{eq:method-start}, first set
\begin{equation}\label{eq:midpoint-stage-one}
 F_{n,1}:=F(t_n,Y_n).
\end{equation}
Apply the reduced \(K\)-\(L\) step to
\((U_n,S_n,V_n)\) with
\begin{equation}\label{eq:midpoint-stage-two-increment}
 \Delta A_{n,2}:=\frac h2F_{n,1}.
\end{equation}
Denote the midpoint approximation and the corresponding value of \(F\) by
\begin{equation}\label{eq:midpoint-stage-two}
 Y_{n,2}:=U_{n,2}S_{n,2}V_{n,2}^\top,
 \qquad F_{n,2}:=F\left(t_n+\frac h2,Y_{n,2}\right).
\end{equation}
Restart from \((U_n,S_n,V_n)\) and apply the same calculation with
\begin{equation}\label{eq:midpoint-final-increment}
 \Delta A_n:=hF_{n,2}.
\end{equation}
Its output is
\begin{equation}\label{eq:midpoint-final-output}
 Y_{n+1}:=U_{n+1}S_{n+1}V_{n+1}^\top.
\end{equation}
Algorithm~\ref{alg:midpoint} gives the factor calculations.

\begin{algorithm}[htbp]
\caption{One step of the projector-splitting explicit midpoint method.}
\label{alg:midpoint}
\begin{algorithmic}[1]
\Require \(t_n,h\), and \(Y_n=U_nS_nV_n^\top\) as in
\eqref{eq:method-start}
\State Set \(F_{n,1}:=F(t_n,Y_n)\).
\State Set \(K_{n,2}:=U_nS_n+\frac h2F_{n,1}V_n\).
\State Compute \(U_{n,2}\widehat S_{n,2}=K_{n,2}\), with
\(U_{n,2}^\top U_{n,2}=I_r\).
\State Set
\(L_{n,2}:=V_nS_n^\top U_n^\top U_{n,2}
+\frac h2F_{n,1}^\top U_{n,2}\).
\State Compute \(V_{n,2}S_{n,2}^\top=L_{n,2}\), with
\(V_{n,2}^\top V_{n,2}=I_r\).
\State Set \(Y_{n,2}:=U_{n,2}S_{n,2}V_{n,2}^\top\) and
\(F_{n,2}:=F(t_n+\frac h2,Y_{n,2})\).
\State Restart from \((U_n,S_n,V_n)\) and set
\(K_{n+1}:=U_nS_n+hF_{n,2}V_n\).
\State Compute \(U_{n+1}\widehat S_{n+1}=K_{n+1}\), with
\(U_{n+1}^\top U_{n+1}=I_r\).
\State Set
\(L_{n+1}:=V_nS_n^\top U_n^\top U_{n+1}
+hF_{n,2}^\top U_{n+1}\).
\State Compute \(V_{n+1}S_{n+1}^\top=L_{n+1}\), with
\(V_{n+1}^\top V_{n+1}=I_r\).
\Ensure \(Y_{n+1}:=U_{n+1}S_{n+1}V_{n+1}^\top\)
\end{algorithmic}
\end{algorithm}
\FloatBarrier

Under the standing assumptions and \eqref{eq:assumption-smoothness}
with \(p=2\), Theorem~\ref{thm:rk-stage-criterion} applies to the midpoint
method for all sufficiently small \(h\) and every rank-admissible
computation.

\begin{proof}
The tableau \eqref{eq:midpoint-tableau} gives
\begin{equation}\label{eq:midpoint-moments}
 \sum_{i=1}^{2}b_i=1,\qquad
 \sum_{i=1}^{2}b_ic_i=\frac12,\qquad b_1=0,\quad b_2=1.
\end{equation}
Lemma~\ref{lem:flow-derivatives} gives the derivative bound in
Theorem~\ref{thm:rk-stage-criterion}.  To verify the stage bound, use
\(F_{n,1}=g_n(0)\) and Taylor's formula:
\begin{equation}\label{eq:midpoint-euler-error}
 \left\|Y_n+\Delta A_{n,2}-A_n\left(\frac h2\right)\right\|
 =\left\|Y_n+\frac h2g_n(0)-A_n\left(\frac h2\right)\right\|
 \le\frac{M_1}{8}h^2.
\end{equation}
Choose \(h_0\le\min(1,h_*)\).  Theorem~\ref{thm:accuracy-transfer},
applied at elapsed time \(h/2\), yields
\begin{align}\label{eq:midpoint-stage-error}
 \left\|Y_{n,2}-A_n\left(\frac h2\right)\right\|
 &\le4\left\|Y_n+\Delta A_{n,2}-A_n\left(\frac h2\right)\right\|
    +\frac{5\sqrt{mn}}2\eps_rh e^{L_Fh/2}\notag\\
 &\le\frac{M_1}{2}h^2+\frac{5\sqrt{mn}}2\eps_rh e^{L_Fh/2}
 \le C(h^2+h\eps_r).
\end{align}
This is the required estimate for the only stage with nonzero final
weight.  All bounds are uniform in \(n\) and the singular values of the
numerical approximations, so Theorem~\ref{thm:rk-stage-criterion} gives the local and global
conclusions with \(p=2\).
\end{proof}

\begin{remark}
The same verification applies to every fixed two-stage explicit
Runge--Kutta method
\[
 \begin{array}{c|cc}
  0&0&0\\
  a&a&0\\ \hline
   &b_1&b_2
 \end{array},
 \qquad 0<a\le1,\qquad b_1+b_2=1,\qquad ab_2=\frac12.
\]
Use
\[
 \Delta A_{n,2}:=ahF_{n,1},\qquad
 F_{n,2}:=F(t_n+ah,Y_{n,2}),\qquad
 \Delta A_n:=h(b_1F_{n,1}+b_2F_{n,2})
\]
in Algorithm~\ref{alg:common-base-rk}.  The first stage is exact,
\(Y_{n,1}=A_n(0)\), and the calculation above gives
\(\|Y_{n,2}-A_n(ah)\|\le C(h^2+h\eps_r)\).
The moment conditions follow from \(b_1+b_2=1\) and \(ab_2=1/2\).
Theorem~\ref{thm:rk-stage-criterion} therefore applies with \(p=2\),
with constants that may depend on the fixed coefficients.
\end{remark}

\subsection{Heun's third-order method}
\label{sec:rk3}

We use the three-stage third-order Runge--Kutta method (RK3)
\begin{equation}\label{eq:heun-tableau}
\begin{array}{c|ccc}
0&0&0&0\\
\frac13&\frac13&0&0\\
\frac23&0&\frac23&0\\ \hline
&\frac14&0&\frac34
\end{array}.
\end{equation}

Starting from \eqref{eq:method-start}, perform the following calculations.
\begin{enumerate}
\item Set \(F_{n,1}:=F(t_n,Y_n)\).  Apply the reduced \(K\)-\(L\) step to
\((U_n,S_n,V_n)\) with
\begin{equation}\label{eq:rk3-stage-two-increment}
 \Delta A_{n,2}:=\frac h3F_{n,1}.
\end{equation}
Denote the stage approximation and the corresponding value of \(F\) by
\begin{equation}\label{eq:rk3-stage-two}
 Y_{n,2}:=U_{n,2}S_{n,2}V_{n,2}^\top,
 \qquad F_{n,2}:=F\left(t_n+\frac h3,Y_{n,2}\right).
\end{equation}

\item Restart from \((U_n,S_n,V_n)\) and apply the same calculation
with
\begin{equation}\label{eq:rk3-stage-three-increment}
 \Delta A_{n,3}:=\frac{2h}{3}F_{n,2}.
\end{equation}
Denote the stage approximation and the corresponding value of \(F\) by
\begin{equation}\label{eq:rk3-stage-three}
 Y_{n,3}:=U_{n,3}S_{n,3}V_{n,3}^\top,
 \qquad F_{n,3}:=F\left(t_n+\frac{2h}{3},Y_{n,3}\right).
\end{equation}

\item Restart once more from \((U_n,S_n,V_n)\), now with
\begin{equation}\label{eq:rk3-final-increment}
 \Delta A_n:=h\left(\frac14F_{n,1}+\frac34F_{n,3}\right).
\end{equation}
The output is \(Y_{n+1}:=U_{n+1}S_{n+1}V_{n+1}^\top\).
\end{enumerate}

Under the standing assumptions and \eqref{eq:assumption-smoothness}
with \(p=3\), Theorem~\ref{thm:rk-stage-criterion} applies to this method
for all sufficiently small \(h\) and every rank-admissible computation.

\begin{proof}
The tableau \eqref{eq:heun-tableau} gives
\begin{equation}\label{eq:heun-moments}
 \sum_{i=1}^{3}b_i=1,\qquad
 \sum_{i=1}^{3}b_ic_i=\frac34\frac23=\frac12,\qquad
 \sum_{i=1}^{3}b_ic_i^2=\frac34\left(\frac23\right)^2=\frac13.
\end{equation}
Since \(b_2=0\) and \(Y_{n,1}=A_n(0)\), the remaining stage condition is
\[
 \left\|Y_{n,3}-A_n\left(\frac{2h}{3}\right)\right\|
 \le C(h^3+h\eps_r).
\]
We obtain it from the second stage.  Taylor's formula gives
\begin{equation}\label{eq:rk3-first-stage-consistency}
 \left\|Y_n+\Delta A_{n,2}-A_n\left(\frac h3\right)\right\|
 =\left\|Y_n+\frac h3g_n(0)-A_n\left(\frac h3\right)\right\|
 \le\frac{M_1}{18}h^2.
\end{equation}
For \(h_0\le\min(1,h_*)\), Theorem~\ref{thm:accuracy-transfer} at elapsed
time \(h/3\) yields
\begin{align}\label{eq:rk3-second-stage-error}
 \left\|Y_{n,2}-A_n\left(\frac h3\right)\right\|
 &\le4\left\|Y_n+\Delta A_{n,2}-A_n\left(\frac h3\right)\right\|
       +\frac{5\sqrt{mn}}3\eps_rh e^{L_Fh/3}\notag\\
 &\le\frac{2M_1}{9}h^2+\frac{5\sqrt{mn}}3\eps_rh e^{L_Fh/3}
 \le C(h^2+h\eps_r).
\end{align}
Next, \eqref{eq:rk3-stage-three-increment}, the Lipschitz bound, and
midpoint quadrature on \([0,2h/3]\) give
\begin{align}
 &\left\|Y_n+\Delta A_{n,3}-A_n\left(\frac{2h}{3}\right)\right\|\notag\\
 &\quad=\left\|Y_n+\frac{2h}{3}F\left(t_n+\frac h3,Y_{n,2}\right)
                  -A_n\left(\frac{2h}{3}\right)\right\|\notag\\
 &\quad\le\frac{2h}{3}L_F
            \left\|Y_{n,2}-A_n\left(\frac h3\right)\right\|
       +\left\|\frac{2h}{3}g_n\left(\frac h3\right)
                    -\int_0^{2h/3}g_n(s)\,ds\right\|\notag\\
 &\quad\le C(h^3+h^2\eps_r)+\frac{M_2}{81}h^3
       \le C(h^3+h^2\eps_r).
 \label{eq:rk3-third-target-error}
\end{align}
Theorem~\ref{thm:accuracy-transfer}, now at elapsed time \(2h/3\), gives
\begin{align}\label{eq:rk3-third-stage-error}
 \left\|Y_{n,3}-A_n\left(\frac{2h}{3}\right)\right\|
 &\le4\left\|Y_n+\Delta A_{n,3}
                -A_n\left(\frac{2h}{3}\right)\right\|
       +\frac{10\sqrt{mn}}3\eps_rh e^{2L_Fh/3}\notag\\
 &\le C(h^3+h\eps_r).
\end{align}
Together with Lemma~\ref{lem:flow-derivatives} and
\eqref{eq:heun-moments}, this verifies every hypothesis of
Theorem~\ref{thm:rk-stage-criterion}, uniformly over the steps.
Its local and global conclusions follow with \(p=3\).
\end{proof}

\begin{remark}
\label{rem:higher-orders}
The classical four-stage fourth-order Runge--Kutta method (RK4) assigns
nonzero final weights to its second and third
stages, whose available bounds are \(O(h^2+h\eps_r)\).  Higher robust
orders can be obtained by adding stages that improve the internal
approximations before they enter the final weighted sum.  For example,
a five-stage construction can use stages two and three as predictors
with \(b_2=b_3=0\), followed by stages at \(c_4=1/2\) and \(c_5=1\)
with errors \(O(h^4+h\eps_r)\).  The final weights
\((b_1,b_4,b_5)=(1/6,2/3,1/6)\) give Simpson quadrature, so
Theorem~\ref{thm:rk-stage-criterion} yields robust order four under uniform
\(C^4\) bounds on \(g_n\) and rank-admissible computations.
\end{remark}

\section{Numerical experiments}
\label{sec:numerical}

We compare the projector-splitting Euler method with the midpoint and
Heun third-order methods of Section~\ref{sec:rk} on a heat equation,
an Allen--Cahn equation, and a linear matrix equation with small singular
values.  For each rank, the numerical and reference solutions start from
the same rank-\(r\) matrix.  We report the relative Frobenius error
\(\|Y_N-A(T)\|/\|A(T)\|\), where \(T=t_0+Nh\).

\subsection{Heat equation}
\label{sec:numerical-heat}

We consider the spatially discretized heat equation from
\cite{CerutiEinkemmerKuschLubich2024},
\begin{equation}\label{eq:numerical-heat}
 \dot A=DA+AD+G,
 \qquad
 G_{ij}:=\sum_{k=1}^{11}10^{1-k}e^{-k(x_i^2+x_j^2)}.
\end{equation}
We use \(m=n=32\) interior grid points on \((-\pi,\pi)\), with
\begin{equation}\label{eq:numerical-heat-grid}
 \Delta x:=\frac{2\pi}{33},\qquad
 x_i:=-\pi+i\Delta x,\qquad
 D:=\frac{1}{\Delta x^2}\operatorname{tridiag}(1,-2,1).
\end{equation}
We impose homogeneous Dirichlet boundary conditions.  Starting with the
entries \(\sin x_i\sin x_j\), we evolve the full matrix equation for
time \(0.1\).  The best rank-\(r\) approximation of this value supplies
the common initial value \(Y_0=A(t_0)\) at \(t_0=0.1\).
We then integrate to \(T=1\), with \(r=4,6,8\) and
\(N=128,256,\ldots,4096\).  Reference values are evaluated by
diagonalizing \(D\).

Figure~\ref{fig:heat-equation}(a) shows convergence at \(r=8\).
The fitted slopes are \(1.00\), \(2.01\), and \(3.01\) for Euler,
midpoint, and Heun RK3, respectively.  Panel (b) shows the effect of rank
on the Heun RK3 error.  At \(r=4\) and \(r=6\), the errors approach
\(2.19\cdot10^{-5}\) and \(2.34\cdot10^{-8}\), respectively, while
the error at \(r=8\) reaches \(2.28\cdot10^{-11}\).
The corresponding best rank-\(r\) approximation errors of the reference
solution at \(T\) are \(2.12\cdot10^{-5}\),
\(2.29\cdot10^{-8}\), and \(2.23\cdot10^{-11}\).

\begin{figure}[htbp]
 \centering
 \includegraphics[width=\textwidth]{experiments/heat_equation.pdf}
 \caption{Relative errors for the heat equation at \(T=1\).
 (a) Euler, midpoint, and Heun RK3 at \(r=8\); slopes are fitted over
 the three largest stepsizes.  (b) Heun RK3 at \(r=4,6,8\).}
 \label{fig:heat-equation}
\end{figure}

\subsection{Allen--Cahn equation}
\label{sec:numerical-allen-cahn}

We next use the matrix Allen--Cahn equation from
\cite{NobileRiffaud2026},
\begin{equation}\label{eq:numerical-allen-cahn}
 \dot A=0.01(DA+AD)+A-A^{\circ3},
\end{equation}
where \(A^{\circ3}\) denotes the entrywise cube.  Here \(m=n=64\),
\(x_i=(i-1)\Delta x\), \(\Delta x=2\pi/63\), and
\(D=(\Delta x)^{-2}\operatorname{tridiag}(1,-2,1)\).
We use the initial function
\begin{equation}\label{eq:numerical-allen-initial}
 u_0(x,y):=
 \frac{(e^{-\tan^2x}+e^{-\tan^2y})\sin x\sin y}
 {1+e^{|\csc(x/2)|}+e^{|\csc(y/2)|}},
 \qquad 0\le x,y\le2\pi,
\end{equation}
with values at singular arguments defined by continuity.
For \(r=4,8,12,16\), the common initial value \(Y_0=A(0)\) is the
best rank-\(r\) approximation of \((u_0(x_i,x_j))_{i,j}\).
We integrate to \(T=10\) with \(N=80,160,\ldots,2560\).
For Heun RK3 at \(r=8,16\), we also include \(N=5120,10240\).

At \(r=16\), the midpoint and Heun RK3 methods give slopes \(1.99\)
and \(3.00\) over the three largest stepsizes; see
Figure~\ref{fig:allen-cahn}(a).  At smaller stepsizes, the Heun RK3 error
is nonmonotone and remains on the scale of \(10^{-9}\).
Panel (b) shows the dependence on rank, including the nonmonotone
small-step errors at \(r=8\).  At the finest stepsizes shown, the
Heun RK3 errors for \(r=4,8,12,16\) are
\(7.77\cdot10^{-4}\), \(1.78\cdot10^{-5}\),
\(2.15\cdot10^{-7}\), and \(3.92\cdot10^{-9}\), respectively.

\begin{figure}[htbp]
 \centering
 \includegraphics[width=\textwidth]{experiments/allen_cahn.pdf}
 \caption{Relative errors for the Allen--Cahn equation at \(T=10\).
 (a) Euler, midpoint, and Heun RK3 at \(r=16\); slopes are fitted over
 the three largest stepsizes.  (b) Heun RK3 at \(r=4,8,12,16\), with
 the additional refinements at \(r=8,16\).}
 \label{fig:allen-cahn}
\end{figure}

\subsection{Small singular values}
\label{sec:numerical-sigma}

Finally, we consider the linear matrix equation
\begin{equation}\label{eq:numerical-linear}
 \dot A=\Omega(t)A+A\Gamma(t),\qquad
 \Omega(t):=\Omega_0+\sin(2t)\Omega_1,\qquad
 \Gamma(t):=\Gamma_0+\cos(3t)\Gamma_1,
\end{equation}
with \(m=n=20\) and \(r=4\).  The four generators are fixed
skew-symmetric matrices with operator norms \(0.65,0.45,0.60,0.40\),
respectively.  This equation preserves singular values.

For fixed orthonormal matrices \(U_0,V_0\), set
\begin{equation}\label{eq:numerical-sigma-initial}
 Y_0=A(0)=U_0\operatorname{diag}(1,10^{-1},10^{-2},\sigma)V_0^\top,
 \qquad \sigma\in\{10^{-2},10^{-6},10^{-10},10^{-14}\}.
\end{equation}
We integrate to \(T=1\) with \(h=2^{-3},\ldots,2^{-7}\).
Figure~\ref{fig:small-sigma}(a) plots the largest relative error over the
four values of \(\sigma\), giving slopes \(1.00\), \(2.00\), and
\(3.00\).  Panel (b) fixes \(h=2^{-7}\) and divides the relative error
by \(h^p\), with \(p=2\) for midpoint and \(p=3\) for Heun RK3.
As \(\sigma\) decreases from \(10^{-2}\) to \(10^{-14}\), these scaled
errors vary by factors of approximately \(1.24\) and \(1.23\),
respectively.

\begin{figure}[htbp]
 \centering
 \includegraphics[width=\textwidth]{experiments/small_singular_values.pdf}
 \caption{Relative errors for the linear matrix equation at \(T=1\).
 (a) Maximum over \(\sigma=10^{-2},10^{-6},10^{-10},10^{-14}\);
 slopes are fitted over the three smallest stepsizes.
 (b) Relative errors divided by \(h^p\) at \(h=2^{-7}\).}
 \label{fig:small-sigma}
\end{figure}
\FloatBarrier

\appendix
\section{Appendix: The rank-preserving comparison curve}
\label{app:comparison-curve}

Under the standing assumptions, we prove
Lemma~\ref{lem:comparison-curve} by first extending the
normal-defect bound from \(\Mr\) to its lower-rank strata, and then constructing
left and right generators whose norms do not contain inverse singular
values.
For this proof, write
\begin{equation}\label{eq:lower-rank-set}
 \Mle:=\{Y\in\R^{m\times n}:\rank(Y)\le r\}.
\end{equation}

\subsection{The normal defect on lower-rank strata}
\label{app:lower-strata}

\begin{lemma}
\label{lem:lower-stratum-defect}
Let \(X\in\Mle\).  If \(x\) and \(y\) are unit vectors satisfying
\begin{equation}\label{eq:lower-normal-directions}
 x\perp\range(X),
 \qquad
 y\perp\range(X^\top),
\end{equation}
then
\begin{equation}\label{eq:lower-stratum-defect}
 |x^\top F(t,X)y|\le\eps_r.
\end{equation}
\end{lemma}

\begin{proof}
For \(\rank(X)=r\), equation \eqref{eq:lower-stratum-defect} follows directly
from \eqref{eq:assumption-normal}.  Let \(q:=\rank(X)<r\) and
\(d:=r-q\).  Since \(r<m,n\),
\begin{align}
 \dim\bigl(\range(X)^\perp\cap x^\perp\bigr)
 &=m-q-1\ge d,                                      \label{eq:left-completion-dimension}\\
 \dim\bigl(\range(X^\top)^\perp\cap y^\perp\bigr)
 &=n-q-1\ge d.                                      \label{eq:right-completion-dimension}
\end{align}
Choose orthonormal vectors \(\widehat u_1,\ldots,\widehat u_d\) in the
first space and \(\widehat v_1,\ldots,\widehat v_d\) in the second, and
define, for \(\tau>0\),
\begin{equation}\label{eq:rank-completion}
 X_\tau:=X+\tau\sum_{\ell=1}^d\widehat u_\ell\widehat v_\ell^\top.
\end{equation}
The new left and right directions are orthogonal to the column and row
spaces of \(X\).  Hence
\begin{equation}\label{eq:rank-completion-properties}
 \rank(X_\tau)=r,
 \qquad
 x\perp\range(X_\tau),
 \qquad
 y\perp\range(X_\tau^\top).
\end{equation}
Equation \eqref{eq:assumption-normal} gives
\(|x^\top F(t,X_\tau)y|\le\eps_r\).  Letting \(\tau\to0\) proves
\eqref{eq:lower-stratum-defect}.
\end{proof}

\subsection{Bounded left and right generators}
\label{app:bounded-generators}

\begin{lemma}
\label{lem:bounded-generators}
For every \(t\in[0,T]\) and \(X\in\Mle\), there are matrices
\(G_L(t,X)\in\R^{m\times m}\), \(G_R(t,X)\in\R^{n\times n}\), and
\(E(t,X)\in\R^{m\times n}\) such that
\begin{equation}\label{eq:generator-decomposition}
 F(t,X)=G_L(t,X)X+XG_R(t,X)+E(t,X),
\end{equation}
with
\begin{equation}\label{eq:generator-bounds}
 \|G_L(t,X)\|_2,\ \|G_R(t,X)\|_2
 \le K_G:=(r+\sqrt r)L_F,
 \qquad
 \|E(t,X)\|\le E_G:=\sqrt{mn}\,\eps_r.
\end{equation}
\end{lemma}

\begin{proof}
Let \(q:=\rank(X)\).  If \(q=0\), then \(X=0\).  Lemma
\ref{lem:lower-stratum-defect}, applied to arbitrary unit vectors \(x,y\),
shows that every entry of \(F(t,0)\) in orthonormal coordinates is bounded
by \(\eps_r\).  Thus
\begin{equation}\label{eq:zero-generator-case}
 G_L(t,0):=0,\qquad G_R(t,0):=0,\qquad E(t,0):=F(t,0)
\end{equation}
satisfies \eqref{eq:generator-decomposition}--\eqref{eq:generator-bounds}.

Suppose \(q\ge1\), and take a compact singular value decomposition (SVD)
\begin{equation}\label{eq:generator-svd}
 X=U\Sigma V^\top
  =\sum_{i=1}^q\sigma_i u_i v_i^\top,
 \qquad \sigma_i>0.
\end{equation}
Define
\begin{align}
 X_i&:=X-\sigma_i u_i v_i^\top,                       \label{eq:delete-one}\\
 X_{ij}&:=X-\sigma_i u_i v_i^\top-\sigma_j u_j v_j^\top
 \quad(i\ne j).                                      \label{eq:delete-two}
\end{align}
Define
\begin{align}
 g_i^L&:=\frac{(I-UU^\top)(F(t,X)-F(t,X_i))v_i}{\sigma_i},
                                                               \label{eq:left-generator-column}\\
 g_i^R&:=\frac{(I-VV^\top)(F(t,X)-F(t,X_i))^\top u_i}{\sigma_i}.
                                                               \label{eq:right-generator-column}
\end{align}
The Lipschitz bound in \eqref{eq:assumption-BL} gives
\begin{equation}\label{eq:normal-generator-bounds}
 \|g_i^L\|_2,\ \|g_i^R\|_2\le L_F.
\end{equation}
With
\(G_\perp^L:=[g_1^L,\ldots,g_q^L]\) and
\(G_\perp^R:=[g_1^R,\ldots,g_q^R]\), it follows that
\begin{equation}\label{eq:normal-generator-matrix-bounds}
 \|G_\perp^L\|_2,\ \|G_\perp^R\|_2\le\sqrt q\,L_F.
\end{equation}

Define \(\Gamma\in\R^{q\times q}\) by
\begin{align}
 \Gamma_{ii}
 &:=\frac{u_i^\top(F(t,X)-F(t,X_i))v_i}{2\sigma_i},
                                                               \label{eq:gamma-diagonal}\\
 \Gamma_{ij}
 &:=\frac{u_i^\top(F(t,X)-F(t,X_{ij}))v_j}
          {\sigma_i+\sigma_j}
 \quad(i\ne j).                                      \label{eq:gamma-offdiagonal}
\end{align}
Since
\(\|X-X_{ij}\|=(\sigma_i^2+\sigma_j^2)^{1/2}
\le\sigma_i+\sigma_j\),
\begin{equation}\label{eq:gamma-bound}
 |\Gamma_{ii}|\le\frac{L_F}{2},\qquad
 |\Gamma_{ij}|\le L_F,\qquad
 \|\Gamma\|_2\le\|\Gamma\|\le qL_F.
\end{equation}
Now set
\begin{equation}\label{eq:generators-defined}
 G_L:=G_\perp^LU^\top+U\Gamma U^\top,
 \qquad
 G_R:=V(G_\perp^R)^\top+V\Gamma V^\top.
\end{equation}
Equations \eqref{eq:normal-generator-matrix-bounds} and
\eqref{eq:gamma-bound} yield
\begin{equation}\label{eq:generators-norm}
 \|G_L\|_2,\ \|G_R\|_2
 \le(\sqrt q+q)L_F\le K_G.
\end{equation}

It remains to bound
\begin{equation}\label{eq:generator-residual-defined}
 E:=F(t,X)-G_LX-XG_R.
\end{equation}
Complete \(U,V\) to orthogonal matrices \([U,U_\perp]\) and
\([V,V_\perp]\).  Direct substitution of
\eqref{eq:left-generator-column}--\eqref{eq:generators-defined} gives
\begin{align}
 (I-UU^\top)EVV^\top
 &=\sum_{i=1}^q (I-UU^\top)F(t,X_i)v_i v_i^\top,       \label{eq:residual-left-normal}\\
 UU^\top E(I-VV^\top)
 &=\sum_{i=1}^q u_i u_i^\top F(t,X_i)(I-VV^\top),      \label{eq:residual-right-normal}\\
 (U^\top EV)_{ii}
 &=u_i^\top F(t,X_i)v_i,                              \label{eq:residual-core-diagonal}\\
 (U^\top EV)_{ij}
 &=u_i^\top F(t,X_{ij})v_j\quad(i\ne j),              \label{eq:residual-core-offdiagonal}\\
 (I-UU^\top)E(I-VV^\top)
 &=(I-UU^\top)F(t,X)(I-VV^\top).                       \label{eq:residual-normal-normal}
\end{align}
For every scalar entry on the right-hand sides, the left and right vectors
are normal directions to \(X_i\), \(X_{ij}\), or \(X\), respectively.
Lemma~\ref{lem:lower-stratum-defect} therefore bounds every entry of
\([U,U_\perp]^\top E[V,V_\perp]\) by \(\eps_r\).  Hence
\begin{equation}\label{eq:generator-residual-bound}
 \|E\|\le\sqrt{mn}\,\eps_r=E_G,
\end{equation}
which completes the proof.
\end{proof}

\subsection{Construction of the comparison curve}
\label{app:generator-flow}

For \(t\in[0,T]\) and \(X\in\Mle\), define
\begin{equation}\label{eq:generator-multifunction}
 \mathcal G(t,X):=
 \left\{(G_L,G_R):
 \begin{array}{l}
  \|G_L\|_2,\ \|G_R\|_2\le K_G,\\
  \|F(t,X)-G_LX-XG_R\|\le E_G
 \end{array}
 \right\}.
\end{equation}
Lemma~\ref{lem:bounded-generators} shows that these sets are nonempty.
They are compact and convex.  The continuity of \(F\) also implies that \(\mathcal G\) has a
closed graph.
The SVD-based generators in Lemma~\ref{lem:bounded-generators} need not vary
continuously when singular values coalesce.  We therefore work with the
entire compact convex set \(\mathcal G(t,X)\) and use a differential
inclusion.

\begin{lemma}
\label{lem:generator-flow}
Let \(Y\in\Mle\), \(0\le t_0\le T-h\), and \(h>0\), and let \(A(s)\)
solve \(\dot A(s)=F(t_0+s,A(s))\), \(A(0)=Y\).  There are
absolutely continuous and invertible matrices
\(\mathcal L:[0,h]\to\R^{m\times m}\) and
\(\mathcal R:[0,h]\to\R^{n\times n}\), with
\(\mathcal L(0)=I_m\) and \(\mathcal R(0)=I_n\), such that
\begin{equation}\label{eq:comparison-curve-defined}
 \widetilde Y(s):=\mathcal L(s)Y\mathcal R(s)
\end{equation}
satisfies, for almost every \(s\in[0,h]\),
\begin{equation}\label{eq:generator-flow-residual}
 \|\dot{\widetilde Y}(s)
       -F(t_0+s,\widetilde Y(s))\|\le E_G.
\end{equation}
Moreover,
\begin{equation}\label{eq:generator-flow-distance}
 \|\widetilde Y(s)-A(s)\|
 \le E_Gs e^{L_Fs}.
\end{equation}
The curve and the right factor satisfy, in addition,
\begin{equation}\label{eq:generator-flow-right-factor}
 \rank(\widetilde Y(s))=\rank(Y),
 \qquad
 \|\mathcal R(s)-I_n\|_2\le e^{K_Gs}-1.
\end{equation}
\end{lemma}

\begin{proof}
On the Euclidean product space
\begin{equation}\label{eq:generator-product-space}
 \mathcal E:=\R^{m\times m}\times\R^{n\times n},
 \qquad
 \|(\mathcal L,\mathcal R)\|_{\mathcal E}^2
 :=\|\mathcal L\|^2+\|\mathcal R\|^2,
\end{equation}
consider the differential inclusion
\begin{equation}\label{eq:generator-inclusion}
 (\dot{\mathcal L},\dot{\mathcal R})
 \in\mathcal H(s,\mathcal L,\mathcal R),
 \qquad
 (\mathcal L(0),\mathcal R(0))=(I_m,I_n),
\end{equation}
where
\begin{equation}\label{eq:generator-inclusion-map}
 \mathcal H(s,\mathcal L,\mathcal R)
 :=\left\{(G_L\mathcal L,\mathcal R G_R):
 (G_L,G_R)\in
 \mathcal G(t_0+s,\mathcal L Y\mathcal R)\right\}.
\end{equation}
The values of \(\mathcal H\) are nonempty, compact, and convex.  Its graph
is closed: for a convergent sequence of states and velocities, the
corresponding generators lie in a fixed compact ball; a convergent
subsequence and the closed graph of \(\mathcal G\) give the limiting
relation in \eqref{eq:generator-inclusion-map}.  The images of bounded sets
are bounded by \eqref{eq:generator-inclusion-growth} below, so the closed
graph implies joint upper semicontinuity in
\((s,\mathcal L,\mathcal R)\).  In particular,
\(s\mapsto\mathcal H(s,\mathcal L,\mathcal R)\) is Borel measurable for
every fixed \((\mathcal L,\mathcal R)\).  Moreover,
\begin{equation}\label{eq:generator-inclusion-growth}
 \sup_{Z\in\mathcal H(s,\mathcal L,\mathcal R)}
 \|Z\|_{\mathcal E}
 \le K_G\|(\mathcal L,\mathcal R)\|_{\mathcal E}.
\end{equation}

Choose
\begin{equation}\label{eq:truncation-radius}
 M_*>e^{K_Gh}\|(I_m,I_n)\|_{\mathcal E}
\end{equation}
and let \(\pi_{M_*}\) be the radial projection of \(\mathcal E\) onto the
closed ball of radius \(M_*\).  Define
\begin{equation}\label{eq:truncated-inclusion-map}
 \widetilde{\mathcal H}(s,\mathcal L,\mathcal R)
 :=\mathcal H\bigl(s,\pi_{M_*}(\mathcal L,\mathcal R)\bigr).
\end{equation}
This map is upper semicontinuous with nonempty compact convex values,
bounded by \(K_GM_*\).  The finite-dimensional existence theorem
\cite[Theorem~5.1]{Deimling1992}, applied on the full space \(\mathcal E\),
gives an absolutely continuous solution of the truncated inclusion on
\([0,h]\).  Before
its first exit from the ball, \eqref{eq:generator-inclusion-growth} and
Gronwall's inequality give
\begin{equation}\label{eq:generator-state-bound}
 \|(\mathcal L(s),\mathcal R(s))\|_{\mathcal E}
 \le e^{K_Gs}\|(I_m,I_n)\|_{\mathcal E}<M_*.
\end{equation}
Thus no first exit occurs, and the solution satisfies
\eqref{eq:generator-inclusion} on \([0,h]\).

The matrices \(\mathcal L\) and \(\mathcal R\) are invertible near
\(s=0\).  On every interval on which they are invertible,
\eqref{eq:generator-inclusion} gives, almost everywhere,
\begin{align}
 \dot{\mathcal L}&=G_L\mathcal L,
 &\frac{d}{ds}\mathcal L^{-1}&=-\mathcal L^{-1}G_L,
 &\|G_L\|_2&\le K_G,                                \label{eq:left-factor-inverse-flow}\\
 \dot{\mathcal R}&=\mathcal R G_R,
 &\frac{d}{ds}\mathcal R^{-1}&=-G_R\mathcal R^{-1},
 &\|G_R\|_2&\le K_G.                                \label{eq:right-factor-inverse-flow}
\end{align}
Gronwall's inequality therefore gives
\begin{equation}\label{eq:factor-flow-inverse-bounds}
 \|\mathcal L(s)\|_2,\ \|\mathcal L(s)^{-1}\|_2,
 \ \|\mathcal R(s)\|_2,\ \|\mathcal R(s)^{-1}\|_2
 \le e^{K_Gs}.
\end{equation}
Consequently neither factor can become singular at a finite first time,
and \eqref{eq:factor-flow-inverse-bounds} holds on \([0,h]\).

By \eqref{eq:generator-inclusion-map}, for almost every \(s\) there is a
pair \((G_L,G_R)\in\mathcal G(t_0+s,\widetilde Y(s))\) such that
\begin{equation}\label{eq:comparison-curve-derivative}
 \dot{\widetilde Y}(s)
 =G_L\widetilde Y(s)+\widetilde Y(s)G_R.
\end{equation}
Equations \eqref{eq:generator-multifunction} and
\eqref{eq:comparison-curve-derivative} prove
\eqref{eq:generator-flow-residual}.  Invertibility of the two factors in
\eqref{eq:comparison-curve-defined} gives the rank identity in
\eqref{eq:generator-flow-right-factor}.  Integrating
\(\dot{\mathcal R}=\mathcal R G_R\) from \(\mathcal R(0)=I_n\), and using
\(\|G_R\|_2\le K_G\) together with
\eqref{eq:factor-flow-inverse-bounds}, gives
\[
 \|\mathcal R(s)-I_n\|_2
 \le\int_0^sK_Ge^{K_G\tau}\,d\tau
 =e^{K_Gs}-1,
\]
which is the second bound in \eqref{eq:generator-flow-right-factor}.

Since \(\widetilde Y(0)=A(0)=Y\), equations
\eqref{eq:assumption-BL} and \eqref{eq:generator-flow-residual} imply
\begin{equation}\label{eq:comparison-distance-integral}
 \|\widetilde Y(s)-A(s)\|
 \le L_F\int_0^s\|\widetilde Y(\tau)-A(\tau)\|\,d\tau+E_Gs.
\end{equation}
Gronwall's inequality proves \eqref{eq:generator-flow-distance}.
\end{proof}

\subsection{Overlap with the old row space}
\label{app:old-row-overlap}

\begin{proof}[Completion of the proof of Lemma~\ref{lem:comparison-curve}]
Apply Lemma~\ref{lem:generator-flow} with \(Y=Y_0\).  Equations
\eqref{eq:generator-flow-residual},
\eqref{eq:generator-flow-distance}, and
\eqref{eq:generator-flow-right-factor} prove
\eqref{eq:comparison-rank}--\eqref{eq:comparison-flow}.

It remains to prove \eqref{eq:comparison-overlap}.  Choose \(h_*>0\) such that
\begin{equation}\label{eq:comparison-hstar}
 e^{K_Gh_*}-1\le\frac12;
\end{equation}
when \(K_G=0\), take \(h_*=1\).  For \(0\le s\le h\le h_*\), set
\begin{equation}\label{eq:old-row-overlap-matrices}
 D_s:=V_0^\top\mathcal R(s)V_0,
 \qquad
 Q_s:=D_s^{-1}V_0^\top\mathcal R(s).
\end{equation}
The right-factor bound in \eqref{eq:generator-flow-right-factor} gives
\begin{equation}\label{eq:old-row-overlap-bounds}
 \|D_s-I_r\|_2\le\frac12,
 \qquad
 \|D_s^{-1}\|_2\le2,
 \qquad
 \|Q_s\|_2\le2e^{K_Gs}\le3.
\end{equation}
Using \(Y_0=Y_0V_0V_0^\top\), we obtain
\begin{align}
 \widetilde Y(s)V_0
 &=\mathcal L(s)Y_0V_0 D_s,                              \label{eq:old-row-product}\\
 (\widetilde Y(s)V_0)Q_s
 &=\mathcal L(s)Y_0V_0V_0^\top\mathcal R(s)
  =\widetilde Y(s).                                  \label{eq:old-row-reconstruction}
\end{align}
Equations \eqref{eq:old-row-overlap-bounds} and
\eqref{eq:old-row-reconstruction} prove
\eqref{eq:comparison-overlap} and complete the proof.
\end{proof}

\section*{Acknowledgement}

This work was partially supported by DOE grant DE-SC0023164, NSF grants DMS-2409858 and IIS-2433957, and DoD MURI grant FA9550-24-1-0254.

\section*{Declaration of AI assistance}
The authors used ChatGPT 5.6 Sol to assist in developing the construction
in Appendix~\ref{app:comparison-curve}, polishing the manuscript, and
performing adversarial checks of the mathematical arguments.  The authors
reviewed and checked all AI-assisted material and take full responsibility
for the correctness and integrity of the final manuscript.



\begingroup
\bibliographystyle{siamplain}
\bibliography{references}
\endgroup

\end{document}
